% MJ81--2C
% Akutes Lungenödem, schwer
% 14.0
% 2013-11-30

:ref:`m-lungenoedem-schwer`
Taktik **Vitale Bedrohung**
Belastung minimieren, rasche
medikamentöse Therpie (veranlassen); ärztliche Therapie
notwendig}

\begin{itemize}
  -   |TxMassVitMK|
  \begin{itemize}
    -   Lagerung: Oberkörper hoch, *wenn möglich Beine
    herabhängen lassen*

    Striktes **Bewegungsverbot!**
    -    Notarzt!
    -   Reanimationsbereitschaft!
  \end{itemize}
  -   Beengende Kleidung öffnen
  -   \Z{Evtl. Gabe von Nitroglycerin-Spray durch einen
  Notfallsanitäter}
  -   Transportentscheidung: Internistische Intensiv- bzw. Überwachungstation
\end{itemize}

\begin{Danger}
  -   Patienten mit Lungenödem und brodelndem Atemgeräusch oder ABCD-Problem 
  sind grundsätzlich notarztpflichtig. 
  
  Auch bei kurzer Transportzeit hat die Stabilisierung des Patienten vor Ort Vorrang.
  
  Bereits der Transport in das Fahrzeug kann gefährlich sein!
\end{Danger}

\Customized{
\begin{description}
  \item[\CustAsboe] Beim kardialen Lungenödem ist die Gabe
  von Nitrolingual Pumpspray gem. Algorithmus durch NFS vorgesehen.
 :cite:`Asb921Memo-AL-2011-000001`\ .
\end{description}
}